% --- Archon physics formalization source begin ---
% archon:physics
% archon:covers PhyXMiniProblems/problem_phyx_mini_0120.lean
% archon:source-report reports/phyx_mini/problem_phyx_mini_0120.source.json

\chapter{Physics problem phyx\_mini\_0120}
\label{ch:PhyXMiniProblems_problem_phyx_mini_0120}

\paragraph{Problem source.}
\#\# Physical scenario

Three-Slit Interference. Monochromatic light with wavelength \$\textbackslash{}lambda\$ is incident on a screen with three narrow slits with separation \$d\$, as shown in figure. Light from the middle slit reaches point \$P\$ with electric field \$E \textbackslash{}cos(\textbackslash{}omega t)\$. From the small-angle approximation, light from the upper and lower slits reaches point \$P\$ with electric fields \$E \textbackslash{}cos(\textbackslash{}omega t + \textbackslash{}phi)\$ and \$E \textbackslash{}cos(\textbackslash{}omega t - \textbackslash{}phi)\$, respectively, where \$\textbackslash{}phi = (2\textbackslash{}pi d \textbackslash{}sin 	heta) / \textbackslash{}lambda\$ is the phase lag and phase lead associated with the different path lengths.

\#\# Question

What is the distance from the central maximum to the closest absolute maximum?

\#\# Answer choices

- A: A: 3.25mm

- B: B: 3.32mm

- C: C: 3.12mm

- D: D: 3.34mm°

\#\# Figure caption (auxiliary; use the image as primary evidence)

The image is a diagram depicting the geometric setup for a diffraction or interference problem, often related to a double-slit experiment. Here's a detailed description:

- Two parallel slits are represented on the left side with a distance "d" between them. This section is highlighted with a light orange shading.
- Three blue lines extend from the slits to a point labeled "P" on the right, indicating the path of waves or light rays. 
- The angles between the middle line and the slanting lines are marked as "θ".
- A vertical line on the right side represents a screen where the interference or diffraction pattern is observed.
- The distances are marked on the diagram: 
  - The horizontal distance from the slits to the point P is labeled "R".
  - The vertical distance on the screen from the center line to point P is labeled "y".
  - The path difference line is labeled "d sin θ".
- The right triangle formed by "d sin θ", the path difference, and "d" helps illustrate the relationship between the slit separation and angle.

This diagram illustrates the relationship between the slits, angles, and distances involved in wave interference.

\#\# Dataset metadata

Wave Optics; Physical Model Grounding Reasoning, Spatial Relation Reasoning

\paragraph{Recorded answer/context.}
A: A: 3.25mm

\paragraph{Figure/image path.}
/root/proposal\_for\_physic/science-mango/phyx\_mini\_run/phyx\_data/test\_image/120.png

\paragraph{Formalization target.}
create a compiling Lean file with sorry bodies at `PhyXMiniProblems/problem\_phyx\_mini\_0120.lean`. The Lean declarations must preserve the physical quantities, dimensions or dimensional roles, figure labels, governing-law hypotheses, and final relation expressed by this problem.
Use Mathlib/Physlib names found through LeanExplore where available. If a physics API is missing, introduce faithful local abstractions rather than scalar placeholder aliases.

\begin{theorem}[Physics formalization target]
\label{thm:physics:phyx_mini_0120:target}
\lean{PhyXMiniProblems.ProblemPhyXMini0120.problem_phyx_mini_0120}
\uses{def:physics:phyx-mini-0120:phyxminiproblems-problemphyxmini0120-lengthquantity, def:physics:phyx-mini-0120:phyxminiproblems-problemphyxmini0120-lengthinmillimeters, def:physics:phyx-mini-0120:phyxminiproblems-problemphyxmini0120-threeslitinterferencesetup, def:physics:phyx-mini-0120:phyxminiproblems-problemphyxmini0120-distancefromcentralinmillimeters, def:physics:phyx-mini-0120:phyxminiproblems-problemphyxmini0120-matchesthreeslitfigure, def:physics:phyx-mini-0120:phyxminiproblems-problemphyxmini0120-hasphysicalparameters, def:physics:phyx-mini-0120:phyxminiproblems-problemphyxmini0120-satisfiesmonochromaticfieldsatp, def:physics:phyx-mini-0120:phyxminiproblems-problemphyxmini0120-satisfiesparaxialthreeslitinterferencelaws, def:physics:phyx-mini-0120:phyxminiproblems-problemphyxmini0120-isclosestpositiveabsolutemaximum, def:physics:phyx-mini-0120:phyxminiproblems-problemphyxmini0120-recordeddatasetanswer, def:physics:phyx-mini-0120:phyxminiproblems-problemphyxmini0120-isuniquematchingdisplayedchoice}
The assigned autoformalize agent should translate this physics problem into Lean declarations in the covered file, with theorem and lemma proof bodies written as `by sorry`.
\end{theorem}
\begin{proof}
This is an autoformalization task, not a proof task. Produce statements that can later be proved by the physics prover without weakening the physical model.
\end{proof}

% --- Archon declaration topology begin ---
\section{Declaration topology}
\begin{definition}[Length Quantity]
  \label{def:physics:phyx-mini-0120:phyxminiproblems-problemphyxmini0120-lengthquantity}
  \lean{PhyXMiniProblems.ProblemPhyXMini0120.LengthQuantity}
  A unit-independent physical quantity carrying the dimension of length
\end{definition}
\begin{proof}
  This is the declared model definition of Length Quantity.
\end{proof}
\begin{definition}[Time Quantity]
  \label{def:physics:phyx-mini-0120:phyxminiproblems-problemphyxmini0120-timequantity}
  \lean{PhyXMiniProblems.ProblemPhyXMini0120.TimeQuantity}
  A unit-independent physical time
\end{definition}
\begin{proof}
  This is the declared model definition of Time Quantity.
\end{proof}
\begin{definition}[Angular Frequency Quantity]
  \label{def:physics:phyx-mini-0120:phyxminiproblems-problemphyxmini0120-angularfrequencyquantity}
  \lean{PhyXMiniProblems.ProblemPhyXMini0120.AngularFrequencyQuantity}
  Angular frequency has inverse-time dimension; radians are dimensionless
\end{definition}
\begin{proof}
  This is the declared model definition of Angular Frequency Quantity.
\end{proof}
\begin{definition}[Electric Field Component Quantity]
  \label{def:physics:phyx-mini-0120:phyxminiproblems-problemphyxmini0120-electricfieldcomponentquantity}
  \lean{PhyXMiniProblems.ProblemPhyXMini0120.ElectricFieldComponentQuantity}
  A scalar electric-field component, with dimension mass * length / (time\textasciicircum{}2 * charge)
\end{definition}
\begin{proof}
  This is the declared model definition of Electric Field Component Quantity.
\end{proof}
\begin{definition}[Irradiance Quantity]
  \label{def:physics:phyx-mini-0120:phyxminiproblems-problemphyxmini0120-irradiancequantity}
  \lean{PhyXMiniProblems.ProblemPhyXMini0120.IrradianceQuantity}
  Irradiance has SI unit W / m\textasciicircum{}2, hence dimension mass / time\textasciicircum{}3
\end{definition}
\begin{proof}
  This is the declared model definition of Irradiance Quantity.
\end{proof}
\begin{definition}[length Readout]
  \label{def:physics:phyx-mini-0120:phyxminiproblems-problemphyxmini0120-lengthreadout}
  \lean{PhyXMiniProblems.ProblemPhyXMini0120.lengthReadout}
  \uses{def:physics:phyx-mini-0120:phyxminiproblems-problemphyxmini0120-lengthquantity}
  Scalar readout of a physical length in a selected length unit
\end{definition}
\begin{proof}
  This is the declared model definition of length Readout.
\end{proof}
\begin{definition}[length In Meters]
  \label{def:physics:phyx-mini-0120:phyxminiproblems-problemphyxmini0120-lengthinmeters}
  \lean{PhyXMiniProblems.ProblemPhyXMini0120.lengthInMeters}
  \uses{def:physics:phyx-mini-0120:phyxminiproblems-problemphyxmini0120-lengthquantity, def:physics:phyx-mini-0120:phyxminiproblems-problemphyxmini0120-lengthreadout}
  Scalar readout of a physical length in meters
\end{definition}
\begin{proof}
  This is the declared model definition of length In Meters.
\end{proof}
\begin{definition}[length In Millimeters]
  \label{def:physics:phyx-mini-0120:phyxminiproblems-problemphyxmini0120-lengthinmillimeters}
  \lean{PhyXMiniProblems.ProblemPhyXMini0120.lengthInMillimeters}
  \uses{def:physics:phyx-mini-0120:phyxminiproblems-problemphyxmini0120-lengthquantity, def:physics:phyx-mini-0120:phyxminiproblems-problemphyxmini0120-lengthreadout}
  Scalar readout of a physical length in millimeters
\end{definition}
\begin{proof}
  This is the declared model definition of length In Millimeters.
\end{proof}
\begin{definition}[time In Seconds]
  \label{def:physics:phyx-mini-0120:phyxminiproblems-problemphyxmini0120-timeinseconds}
  \lean{PhyXMiniProblems.ProblemPhyXMini0120.timeInSeconds}
  \uses{def:physics:phyx-mini-0120:phyxminiproblems-problemphyxmini0120-timequantity}
  Scalar SI readout of a physical time, in seconds
\end{definition}
\begin{proof}
  This is the declared model definition of time In Seconds.
\end{proof}
\begin{definition}[angular Frequency In Radians Per Second]
  \label{def:physics:phyx-mini-0120:phyxminiproblems-problemphyxmini0120-angularfrequencyinradianspersecond}
  \lean{PhyXMiniProblems.ProblemPhyXMini0120.angularFrequencyInRadiansPerSecond}
  \uses{def:physics:phyx-mini-0120:phyxminiproblems-problemphyxmini0120-angularfrequencyquantity}
  Scalar SI readout of angular frequency, in radians per second
\end{definition}
\begin{proof}
  This is the declared model definition of angular Frequency In Radians Per Second.
\end{proof}
\begin{definition}[electric Field Component In SI]
  \label{def:physics:phyx-mini-0120:phyxminiproblems-problemphyxmini0120-electricfieldcomponentinsi}
  \lean{PhyXMiniProblems.ProblemPhyXMini0120.electricFieldComponentInSI}
  \uses{def:physics:phyx-mini-0120:phyxminiproblems-problemphyxmini0120-electricfieldcomponentquantity}
  Scalar SI readout of the electric-field component, in volts per meter
\end{definition}
\begin{proof}
  This is the declared model definition of electric Field Component In SI.
\end{proof}
\begin{definition}[irradiance In SI]
  \label{def:physics:phyx-mini-0120:phyxminiproblems-problemphyxmini0120-irradianceinsi}
  \lean{PhyXMiniProblems.ProblemPhyXMini0120.irradianceInSI}
  \uses{def:physics:phyx-mini-0120:phyxminiproblems-problemphyxmini0120-irradiancequantity}
  Scalar SI irradiance readout, in watts per square meter
\end{definition}
\begin{proof}
  This is the declared model definition of irradiance In SI.
\end{proof}
\begin{definition}[Slit Label]
  \label{def:physics:phyx-mini-0120:phyxminiproblems-problemphyxmini0120-slitlabel}
  \lean{PhyXMiniProblems.ProblemPhyXMini0120.SlitLabel}
  The three apertures, ordered vertically as in the source image
\end{definition}
\begin{proof}
  This is the declared model definition of Slit Label.
\end{proof}
\begin{definition}[Optical Plane]
  \label{def:physics:phyx-mini-0120:phyxminiproblems-problemphyxmini0120-opticalplane}
  \lean{PhyXMiniProblems.ProblemPhyXMini0120.OpticalPlane}
  The two optical planes shown in the ray diagram
\end{definition}
\begin{proof}
  This is the declared model definition of Optical Plane.
\end{proof}
\begin{definition}[Screen Point]
  \label{def:physics:phyx-mini-0120:phyxminiproblems-problemphyxmini0120-screenpoint}
  \lean{PhyXMiniProblems.ProblemPhyXMini0120.ScreenPoint}
  The central screen point and the explicitly labeled point P
\end{definition}
\begin{proof}
  This is the declared model definition of Screen Point.
\end{proof}
\begin{definition}[Figure Axis]
  \label{def:physics:phyx-mini-0120:phyxminiproblems-problemphyxmini0120-figureaxis}
  \lean{PhyXMiniProblems.ProblemPhyXMini0120.FigureAxis}
  Coordinate axes implicit in the labels R (axial) and y (transverse)
\end{definition}
\begin{proof}
  This is the declared model definition of Figure Axis.
\end{proof}
\begin{definition}[Aperture Model]
  \label{def:physics:phyx-mini-0120:phyxminiproblems-problemphyxmini0120-aperturemodel}
  \lean{PhyXMiniProblems.ProblemPhyXMini0120.ApertureModel}
  Whether an aperture is idealized as narrow or assigned a finite width
\end{definition}
\begin{proof}
  This is the declared model definition of Aperture Model.
\end{proof}
\begin{definition}[Optical Approximation]
  \label{def:physics:phyx-mini-0120:phyxminiproblems-problemphyxmini0120-opticalapproximation}
  \lean{PhyXMiniProblems.ProblemPhyXMini0120.OpticalApproximation}
  The propagation model used to project a ray angle onto the screen
\end{definition}
\begin{proof}
  This is the declared model definition of Optical Approximation.
\end{proof}
\begin{definition}[Three Slit Interference Setup]
  \label{def:physics:phyx-mini-0120:phyxminiproblems-problemphyxmini0120-threeslitinterferencesetup}
  \lean{PhyXMiniProblems.ProblemPhyXMini0120.ThreeSlitInterferenceSetup}
  \uses{def:physics:phyx-mini-0120:phyxminiproblems-problemphyxmini0120-lengthquantity, def:physics:phyx-mini-0120:phyxminiproblems-problemphyxmini0120-timequantity, def:physics:phyx-mini-0120:phyxminiproblems-problemphyxmini0120-angularfrequencyquantity, def:physics:phyx-mini-0120:phyxminiproblems-problemphyxmini0120-electricfieldcomponentquantity, def:physics:phyx-mini-0120:phyxminiproblems-problemphyxmini0120-irradiancequantity, def:physics:phyx-mini-0120:phyxminiproblems-problemphyxmini0120-slitlabel, def:physics:phyx-mini-0120:phyxminiproblems-problemphyxmini0120-opticalplane, def:physics:phyx-mini-0120:phyxminiproblems-problemphyxmini0120-screenpoint, def:physics:phyx-mini-0120:phyxminiproblems-problemphyxmini0120-figureaxis, def:physics:phyx-mini-0120:phyxminiproblems-problemphyxmini0120-aperturemodel, def:physics:phyx-mini-0120:phyxminiproblems-problemphyxmini0120-opticalapproximation}
  Physical quantities, field readouts, and labeled geometry for the experiment. phaseAtOffsetY y is the adjacent-slit phase at signed screen offset y. The three individual fields at P and their coherent sum remain dimensionful electric-field components. The full screen pattern is a dimensionful irradiance
\end{definition}
\begin{proof}
  This is the declared model definition of Three Slit Interference Setup.
\end{proof}
\begin{definition}[slit Phase At P]
  \label{def:physics:phyx-mini-0120:phyxminiproblems-problemphyxmini0120-slitphaseatp}
  \lean{PhyXMiniProblems.ProblemPhyXMini0120.slitPhaseAtP}
  \uses{def:physics:phyx-mini-0120:phyxminiproblems-problemphyxmini0120-slitlabel, def:physics:phyx-mini-0120:phyxminiproblems-problemphyxmini0120-threeslitinterferencesetup}
  The phase assigned to each slit in the field formulas stated at P
\end{definition}
\begin{proof}
  This is the declared model definition of slit Phase At P.
\end{proof}
\begin{definition}[signed Offset From Central]
  \label{def:physics:phyx-mini-0120:phyxminiproblems-problemphyxmini0120-signedoffsetfromcentral}
  \lean{PhyXMiniProblems.ProblemPhyXMini0120.signedOffsetFromCentral}
  \uses{def:physics:phyx-mini-0120:phyxminiproblems-problemphyxmini0120-lengthquantity, def:physics:phyx-mini-0120:phyxminiproblems-problemphyxmini0120-lengthreadout, def:physics:phyx-mini-0120:phyxminiproblems-problemphyxmini0120-threeslitinterferencesetup}
  Signed transverse separation from the central screen point, in a unit
\end{definition}
\begin{proof}
  This is the declared model definition of signed Offset From Central.
\end{proof}
\begin{definition}[distance From Central In Millimeters]
  \label{def:physics:phyx-mini-0120:phyxminiproblems-problemphyxmini0120-distancefromcentralinmillimeters}
  \lean{PhyXMiniProblems.ProblemPhyXMini0120.distanceFromCentralInMillimeters}
  \uses{def:physics:phyx-mini-0120:phyxminiproblems-problemphyxmini0120-lengthquantity, def:physics:phyx-mini-0120:phyxminiproblems-problemphyxmini0120-threeslitinterferencesetup, def:physics:phyx-mini-0120:phyxminiproblems-problemphyxmini0120-signedoffsetfromcentral}
  Positive distance from the central point, read in millimeters
\end{definition}
\begin{proof}
  This is the declared model definition of distance From Central In Millimeters.
\end{proof}
\begin{definition}[Matches Three Slit Figure]
  \label{def:physics:phyx-mini-0120:phyxminiproblems-problemphyxmini0120-matchesthreeslitfigure}
  \lean{PhyXMiniProblems.ProblemPhyXMini0120.MatchesThreeSlitFigure}
  \uses{def:physics:phyx-mini-0120:phyxminiproblems-problemphyxmini0120-lengthreadout, def:physics:phyx-mini-0120:phyxminiproblems-problemphyxmini0120-lengthinmeters, def:physics:phyx-mini-0120:phyxminiproblems-problemphyxmini0120-slitlabel, def:physics:phyx-mini-0120:phyxminiproblems-problemphyxmini0120-threeslitinterferencesetup}
  Qualitative and metric information read from the primary image. The upper-to-middle and middle-to-lower separations are both d; every ray arrives at P; the screen is R beyond the slit plane; the central point has zero transverse offset; and the signed adjacent path differences at P are d sin theta, 0, and -d sin theta. No requested fringe distance occurs in this figure predicate
\end{definition}
\begin{proof}
  This is the declared model definition of Matches Three Slit Figure.
\end{proof}
\begin{definition}[Has Physical Parameters]
  \label{def:physics:phyx-mini-0120:phyxminiproblems-problemphyxmini0120-hasphysicalparameters}
  \lean{PhyXMiniProblems.ProblemPhyXMini0120.HasPhysicalParameters}
  \uses{def:physics:phyx-mini-0120:phyxminiproblems-problemphyxmini0120-lengthinmeters, def:physics:phyx-mini-0120:phyxminiproblems-problemphyxmini0120-angularfrequencyinradianspersecond, def:physics:phyx-mini-0120:phyxminiproblems-problemphyxmini0120-electricfieldcomponentinsi, def:physics:phyx-mini-0120:phyxminiproblems-problemphyxmini0120-irradianceinsi, def:physics:phyx-mini-0120:phyxminiproblems-problemphyxmini0120-threeslitinterferencesetup}
  Positivity conditions for the physical inputs and central illumination
\end{definition}
\begin{proof}
  This is the declared model definition of Has Physical Parameters.
\end{proof}
\begin{definition}[Satisfies Monochromatic Fields At P]
  \label{def:physics:phyx-mini-0120:phyxminiproblems-problemphyxmini0120-satisfiesmonochromaticfieldsatp}
  \lean{PhyXMiniProblems.ProblemPhyXMini0120.SatisfiesMonochromaticFieldsAtP}
  \uses{def:physics:phyx-mini-0120:phyxminiproblems-problemphyxmini0120-timequantity, def:physics:phyx-mini-0120:phyxminiproblems-problemphyxmini0120-lengthinmeters, def:physics:phyx-mini-0120:phyxminiproblems-problemphyxmini0120-timeinseconds, def:physics:phyx-mini-0120:phyxminiproblems-problemphyxmini0120-angularfrequencyinradianspersecond, def:physics:phyx-mini-0120:phyxminiproblems-problemphyxmini0120-electricfieldcomponentinsi, def:physics:phyx-mini-0120:phyxminiproblems-problemphyxmini0120-slitlabel, def:physics:phyx-mini-0120:phyxminiproblems-problemphyxmini0120-threeslitinterferencesetup, def:physics:phyx-mini-0120:phyxminiproblems-problemphyxmini0120-slitphaseatp}
  The three monochromatic field formulas stated in the problem, together with linear coherent superposition. The phase convention is upper +phi, middle 0, and lower -phi
\end{definition}
\begin{proof}
  This is the declared model definition of Satisfies Monochromatic Fields At P.
\end{proof}
\begin{definition}[Satisfies Paraxial Three Slit Interference Laws]
  \label{def:physics:phyx-mini-0120:phyxminiproblems-problemphyxmini0120-satisfiesparaxialthreeslitinterferencelaws}
  \lean{PhyXMiniProblems.ProblemPhyXMini0120.SatisfiesParaxialThreeSlitInterferenceLaws}
  \uses{def:physics:phyx-mini-0120:phyxminiproblems-problemphyxmini0120-lengthquantity, def:physics:phyx-mini-0120:phyxminiproblems-problemphyxmini0120-lengthinmeters, def:physics:phyx-mini-0120:phyxminiproblems-problemphyxmini0120-irradianceinsi, def:physics:phyx-mini-0120:phyxminiproblems-problemphyxmini0120-threeslitinterferencesetup, def:physics:phyx-mini-0120:phyxminiproblems-problemphyxmini0120-signedoffsetfromcentral}
  Governing small-angle geometry and ideal equal-amplitude three-slit interference across the screen. The laws are y = R sin theta, phi = 2 pi d sin(theta) / lambda, and the normalized intensity I = I\_central ((1 + 2 cos phi) / 3)\textasciicircum{}2. They contain neither the requested 3.25 mm value nor an answer-choice label
\end{definition}
\begin{proof}
  This is the declared model definition of Satisfies Paraxial Three Slit Interference Laws.
\end{proof}
\begin{definition}[Is Absolute Intensity Maximum]
  \label{def:physics:phyx-mini-0120:phyxminiproblems-problemphyxmini0120-isabsoluteintensitymaximum}
  \lean{PhyXMiniProblems.ProblemPhyXMini0120.IsAbsoluteIntensityMaximum}
  \uses{def:physics:phyx-mini-0120:phyxminiproblems-problemphyxmini0120-lengthquantity, def:physics:phyx-mini-0120:phyxminiproblems-problemphyxmini0120-irradianceinsi, def:physics:phyx-mini-0120:phyxminiproblems-problemphyxmini0120-threeslitinterferencesetup}
  A screen offset realizes the global (absolute) maximum irradiance
\end{definition}
\begin{proof}
  This is the declared model definition of Is Absolute Intensity Maximum.
\end{proof}
\begin{definition}[Is Closest Positive Absolute Maximum]
  \label{def:physics:phyx-mini-0120:phyxminiproblems-problemphyxmini0120-isclosestpositiveabsolutemaximum}
  \lean{PhyXMiniProblems.ProblemPhyXMini0120.IsClosestPositiveAbsoluteMaximum}
  \uses{def:physics:phyx-mini-0120:phyxminiproblems-problemphyxmini0120-lengthquantity, def:physics:phyx-mini-0120:phyxminiproblems-problemphyxmini0120-threeslitinterferencesetup, def:physics:phyx-mini-0120:phyxminiproblems-problemphyxmini0120-signedoffsetfromcentral, def:physics:phyx-mini-0120:phyxminiproblems-problemphyxmini0120-isabsoluteintensitymaximum}
  The closest absolute maximum on the positive-y side of the central maximum. This identifies the requested point without assigning its distance
\end{definition}
\begin{proof}
  This is the declared model definition of Is Closest Positive Absolute Maximum.
\end{proof}
\begin{definition}[Answer Choice]
  \label{def:physics:phyx-mini-0120:phyxminiproblems-problemphyxmini0120-answerchoice}
  \lean{PhyXMiniProblems.ProblemPhyXMini0120.AnswerChoice}
  The four answer labels printed with the problem
\end{definition}
\begin{proof}
  This is the declared model definition of Answer Choice.
\end{proof}
\begin{definition}[displayed Distance Millimeters]
  \label{def:physics:phyx-mini-0120:phyxminiproblems-problemphyxmini0120-displayeddistancemillimeters}
  \lean{PhyXMiniProblems.ProblemPhyXMini0120.displayedDistanceMillimeters}
  \uses{def:physics:phyx-mini-0120:phyxminiproblems-problemphyxmini0120-answerchoice}
  Displayed distance readouts, in millimeters. The degree sign appended to D in the source is dimensionally inconsistent and is not part of its readout
\end{definition}
\begin{proof}
  This is the declared model definition of displayed Distance Millimeters.
\end{proof}
\begin{definition}[recorded Dataset Answer]
  \label{def:physics:phyx-mini-0120:phyxminiproblems-problemphyxmini0120-recordeddatasetanswer}
  \lean{PhyXMiniProblems.ProblemPhyXMini0120.recordedDatasetAnswer}
  \uses{def:physics:phyx-mini-0120:phyxminiproblems-problemphyxmini0120-answerchoice}
  The answer label recorded in the source dataset
\end{definition}
\begin{proof}
  This is the declared model definition of recorded Dataset Answer.
\end{proof}
\begin{definition}[Matches Displayed Distance Choice]
  \label{def:physics:phyx-mini-0120:phyxminiproblems-problemphyxmini0120-matchesdisplayeddistancechoice}
  \lean{PhyXMiniProblems.ProblemPhyXMini0120.MatchesDisplayedDistanceChoice}
  \uses{def:physics:phyx-mini-0120:phyxminiproblems-problemphyxmini0120-lengthquantity, def:physics:phyx-mini-0120:phyxminiproblems-problemphyxmini0120-threeslitinterferencesetup, def:physics:phyx-mini-0120:phyxminiproblems-problemphyxmini0120-distancefromcentralinmillimeters, def:physics:phyx-mini-0120:phyxminiproblems-problemphyxmini0120-answerchoice, def:physics:phyx-mini-0120:phyxminiproblems-problemphyxmini0120-displayeddistancemillimeters}
  A displayed choice equals the requested central-to-maximum distance
\end{definition}
\begin{proof}
  This is the declared model definition of Matches Displayed Distance Choice.
\end{proof}
\begin{definition}[Is Unique Matching Displayed Choice]
  \label{def:physics:phyx-mini-0120:phyxminiproblems-problemphyxmini0120-isuniquematchingdisplayedchoice}
  \lean{PhyXMiniProblems.ProblemPhyXMini0120.IsUniqueMatchingDisplayedChoice}
  \uses{def:physics:phyx-mini-0120:phyxminiproblems-problemphyxmini0120-lengthquantity, def:physics:phyx-mini-0120:phyxminiproblems-problemphyxmini0120-threeslitinterferencesetup, def:physics:phyx-mini-0120:phyxminiproblems-problemphyxmini0120-answerchoice, def:physics:phyx-mini-0120:phyxminiproblems-problemphyxmini0120-matchesdisplayeddistancechoice}
  A choice is the unique displayed distance matching the requested point
\end{definition}
\begin{proof}
  This is the declared model definition of Is Unique Matching Displayed Choice.
\end{proof}
\begin{lemma}[resultant Electric Field At P]
  \label{lem:physics:phyx-mini-0120:phyxminiproblems-problemphyxmini0120-resultantelectricfieldatp}
  \lean{PhyXMiniProblems.ProblemPhyXMini0120.resultantElectricFieldAtP}
  \uses{def:physics:phyx-mini-0120:phyxminiproblems-problemphyxmini0120-timequantity, def:physics:phyx-mini-0120:phyxminiproblems-problemphyxmini0120-timeinseconds, def:physics:phyx-mini-0120:phyxminiproblems-problemphyxmini0120-angularfrequencyinradianspersecond, def:physics:phyx-mini-0120:phyxminiproblems-problemphyxmini0120-electricfieldcomponentinsi, def:physics:phyx-mini-0120:phyxminiproblems-problemphyxmini0120-threeslitinterferencesetup, def:physics:phyx-mini-0120:phyxminiproblems-problemphyxmini0120-satisfiesmonochromaticfieldsatp}
  The stated harmonic components coherently add to E (1 + 2 cos phi) cos(omega t) at P
\end{lemma}
\begin{proof}
  The conclusion follows from the governing relations and figure readouts named in the statement of resultant Electric Field At P.
\end{proof}
\begin{lemma}[closest Absolute Maximum phase eq two pi]
  \label{lem:physics:phyx-mini-0120:phyxminiproblems-problemphyxmini0120-closestabsolutemaximum-phase-eq-two-pi}
  \lean{PhyXMiniProblems.ProblemPhyXMini0120.closestAbsoluteMaximum_phase_eq_two_pi}
  \uses{def:physics:phyx-mini-0120:phyxminiproblems-problemphyxmini0120-lengthquantity, def:physics:phyx-mini-0120:phyxminiproblems-problemphyxmini0120-threeslitinterferencesetup, def:physics:phyx-mini-0120:phyxminiproblems-problemphyxmini0120-matchesthreeslitfigure, def:physics:phyx-mini-0120:phyxminiproblems-problemphyxmini0120-hasphysicalparameters, def:physics:phyx-mini-0120:phyxminiproblems-problemphyxmini0120-satisfiesparaxialthreeslitinterferencelaws, def:physics:phyx-mini-0120:phyxminiproblems-problemphyxmini0120-isclosestpositiveabsolutemaximum}
  In the paraxial model the first positive absolute maximum has adjacent-slit phase 2 pi, the first recurrence of the central phasor alignment
\end{lemma}
\begin{proof}
  The conclusion follows from the governing relations and figure readouts named in the statement of closest Absolute Maximum phase eq two pi.
\end{proof}
\begin{lemma}[closest Absolute Maximum Distance formula]
  \label{lem:physics:phyx-mini-0120:phyxminiproblems-problemphyxmini0120-closestabsolutemaximumdistance-formula}
  \lean{PhyXMiniProblems.ProblemPhyXMini0120.closestAbsoluteMaximumDistance_formula}
  \uses{def:physics:phyx-mini-0120:phyxminiproblems-problemphyxmini0120-lengthquantity, def:physics:phyx-mini-0120:phyxminiproblems-problemphyxmini0120-lengthinmillimeters, def:physics:phyx-mini-0120:phyxminiproblems-problemphyxmini0120-threeslitinterferencesetup, def:physics:phyx-mini-0120:phyxminiproblems-problemphyxmini0120-distancefromcentralinmillimeters, def:physics:phyx-mini-0120:phyxminiproblems-problemphyxmini0120-matchesthreeslitfigure, def:physics:phyx-mini-0120:phyxminiproblems-problemphyxmini0120-hasphysicalparameters, def:physics:phyx-mini-0120:phyxminiproblems-problemphyxmini0120-satisfiesparaxialthreeslitinterferencelaws, def:physics:phyx-mini-0120:phyxminiproblems-problemphyxmini0120-isclosestpositiveabsolutemaximum}
  The symbolic distance to the closest positive absolute maximum is R lambda / d, expressed using millimeter readouts throughout
\end{lemma}
\begin{proof}
  The conclusion follows from the governing relations and figure readouts named in the statement of closest Absolute Maximum Distance formula.
\end{proof}
% --- Archon declaration topology end ---
% --- Archon physics formalization source end ---
